\begin{textsample}{POS Dim 4 – human – Score 6.00 – es/es\_ef\_anos\_iniciais\_linguagens\_lp\_2\_p7.txt}  \label{ex:f4_pos_005}
No \textbf{eixo} Leitura/Escuta, o uso social da \textbf{leitura} e da \textbf{escrita} se efetiva por meio da progressiva incorporação de estratégias de \textbf{leitura} em \textbf{textos} de nível de complexidade crescente, assim como no \textbf{eixo} Produção de \textbf{textos}, pela incorporação de estratégias de diferentes gêneros \textbf{textuais}, de forma também progressiva.
Os \textbf{eixos} do componente curricular de Língua Portuguesa estão mutuamente \textbf{associados} entre si.
Para participar com autonomia dos diferentes contextos, o estudante \textbf{precisa} falar, escutar, ler e escrever, compreendendo, interpretando, produzindo, observando os aspectos culturais da língua materna, os conhecimentos linguísticos e gramaticais, a estrutura dos gêneros \textbf{textuais} literários e não literários.
Tendo esse fato como pano de fundo, cumpre discorrer sobre os conceitos e as definições dos \textbf{eixos} presentes neste documento.

% matched lemmas: associar, eixo, escrita, leitura, precisar, texto, textual
\end{textsample}
